%% results_detailed.tex — TWINGATE Supplementary Results
%% Companion to main.tex. Standalone document (article class).
%% Purpose: full derivation detail for every number that appears in
%% main.tex (condensed there for narrative flow), plus explored
%% directions not currently included in main.tex, plus the provenance
%% history of numbers that were corrected during this project.
%% Use this document as a menu: promote any section's content into
%% main.tex if a reviewer or future draft calls for more depth there.

\documentclass[11pt]{article}
\usepackage[margin=1in]{geometry}
\usepackage{amsmath,amssymb}
\usepackage{booktabs}
\usepackage{longtable}
\usepackage{hyperref}
\usepackage{xcolor}
\usepackage{parskip}

\title{TWINGATE: Supplementary Results and Derivation Detail\\[6pt]
\large Companion document to \emph{TWINGATE: A Twin-Gated Framework for
Physics-Based Coverage Validation in V2X Digital Twins}}
\author{}
\date{}

\begin{document}
\maketitle

\tableofcontents

\section{How to use this document}

Everything in \texttt{main.tex} is condensed for narrative flow: headline
numbers, the theoretical reasoning behind them, and pointers to ``full
detail here.'' This document \emph{is} that full detail. Each section
below is self-contained and mapped to the main-paper section it supports,
so material can be promoted back into the main paper (e.g., if a reviewer
asks for it) without re-deriving anything.

Section~\ref{sec:m4} and Section~\ref{sec:handover} describe experiments
that were run, produced real results, and are \textbf{not currently
included in main.tex} at all — they are genuine findings (one positive,
one a well-diagnosed negative result) held in reserve pending a decision
on whether they strengthen the paper enough to include.

\section{Pipeline Hyperparameters}
\label{sec:hyperparams}

Supports: Section III.A (3D Scene Construction), III.D.3 (4D Optimization).

\subsection{Scene construction}
\begin{itemize}
  \item Operator~1 scene origin: $(52.507005^\circ\text{N}, 13.323428^\circ\text{E})$
  \item Operator~2 scene origin: $(52.506112^\circ\text{N}, 13.321908^\circ\text{E})$
  \item Ray tracing backend: Sionna RT 2.0.1 on the \texttt{cuda\_ad\_mono\_polarized}
        Mitsuba~3 variant (GPU-accelerated, monochromatic, polarized).
  \item Max path depth: 5 interactions (reflection + diffraction; scattering
        added in Stage~4, see below).
  \item Per-call cost: a batched forward pass over 40 receiver locations
        takes approximately 2--3\,s on an NVIDIA GPU under this configuration.
        A full 4D Nelder-Mead optimization of one tower requires
        80--95 such calls (including the coarse azimuth grid), i.e.\
        roughly 3--6 minutes per tower on a single GPU.
\end{itemize}

\subsection{4D Nelder-Mead optimizer (Gate 2 position/height/azimuth search)}
\begin{itemize}
  \item Search variables: $(\Delta x, \Delta y, h, \varphi)$ — horizontal
        displacement from WCL init, height $h \in [10,60]$\,m, azimuth
        $\varphi \in [0^\circ,360^\circ)$.
  \item Coarse pre-search: 8-direction azimuth grid at
        $0^\circ, 45^\circ, 90^\circ, \ldots, 315^\circ$; best-scoring orientation seeds the
        full 4D simplex.
  \item Simplex initialization: 5 vertices at the WCL position, with a
        50\,m horizontal arm, 10\,m height arm, and $45^\circ$ azimuth arm.
  \item Sionna yaw convention: compass bearing $\varphi$ maps to scene yaw
        via $\alpha = \pi/2 - \varphi$ (x-axis = East).
  \item Search bound: 500\,m radius in $(\Delta x, \Delta y)$; height and
        azimuth clamped to their physical ranges at every evaluation.
  \item Convergence: $|\Delta\boldsymbol\theta| < 10$\,m and
        $|\Delta\text{MAE}| < 0.02$\,dB, 80-iteration hard cap.
  \item Per-objective-evaluation training subsample: 40 measurement rows
        (one GPU batch) during the search; the final evaluation pass uses
        all available training and validation rows for that tower.
  \item OLS calibration slope bound: $\alpha_i \in [0, 2]$ (prevents sign
        flips and unphysical amplification).
\end{itemize}

\subsection{Diffuse scattering (Stage 4)}
\begin{itemize}
  \item Material: ITU \texttt{concrete} preset (all building surfaces).
  \item Scattering coefficient: $S = 0.4$ (TR~38.901 value at 2\,GHz),
        vs.\ Sionna's default $S=0$.
  \item Scattering pattern: Lambertian (isotropic diffuse lobe).
  \item Everything else in the pipeline (NM search, sector antenna, OLS
        calibration) is held fixed between Stage~3 ($S=0$) and Stage~4
        ($S=0.4$) to isolate the marginal contribution of scattering.
\end{itemize}

\section{Per-Device Gate 2 Results (Full Breakdown)}
\label{sec:perdevice}

Supports: Table~I and \S IV.C (Gate 2: RSRP Fidelity) in main.tex, which
report only the Operator-level combined figures.

\begin{table}[h]
\centering
\caption{Per-device Stage 3 (position/orientation refinement, no
scattering) results, temporal split, held-out day 3.}
\small
\begin{tabular}{lccccc}
\toprule
Device & Towers & WCL MAE (dB) & Refined MAE (dB) & Gain (dB) & Improved \\
\midrule
pc1 & 81  & 3.877 & 3.029 & +0.848 & 63/81 (77.8\%) \\
pc4 & 43\textsuperscript{a} & 3.787 & 3.016 & +0.771 & 33/43 (76.7\%) \\
pc2 & 37  & 6.927 & 4.601 & +2.327 & 25/37 (67.6\%) \\
pc3 & 67  & 6.786 & 3.986 & +2.801 & 53/67 (79.1\%) \\
\midrule
Op.\,1 (pc1+pc4) & 124 & 3.865 & 3.027 & +0.838 & 96/124 (77.4\%) \\
Op.\,2 (pc2+pc3) & 104 & 6.805 & 4.069 & +2.736 & 78/104 (75.0\%) \\
\bottomrule
\end{tabular}\\[2mm]
\raggedright\footnotesize\textsuperscript{a}45 towers refined; 2 excluded
(see \S\ref{sec:nan}).
\end{table}

\begin{table}[h]
\centering
\caption{Per-device Stage 4 (+ diffuse scattering, $S=0.4$) results,
temporal split, held-out day 3.}
\small
\begin{tabular}{lccccc}
\toprule
Device & Towers & WCL MAE (dB) & Refined MAE (dB) & Gain (dB) & Improved \\
\midrule
pc1 & 81  & 3.874 & 3.020 & +0.853 & 62/81 (76.5\%) \\
pc4 & 43\textsuperscript{a} & 3.803 & 2.980 & +0.823 & 32/43 (74.4\%) \\
pc2 & 37  & 6.939 & 4.553 & +2.387 & 26/37 (70.3\%) \\
pc3 & 67  & 6.743 & 3.957 & +2.786 & 54/67 (80.6\%) \\
\midrule
Op.\,1 (pc1+pc4) & 124 & 3.865 & 3.015 & +0.849 & 94/124 (75.8\%) \\
Op.\,2 (pc2+pc3) & 104 & 6.770 & 4.037 & +2.732 & 80/104 (76.9\%) \\
\bottomrule
\end{tabular}
\end{table}

Median/mean TX position correction (Stage~4, combined Op.\,1):
median 45\,m, mean 63\,m. Refined TX heights span the 10--60\,m search
range with recovered azimuths spanning the full compass, i.e.\ no
boundary-clamping artifact.

Total Sionna calls per run (full pipeline, one device, one stage):
pc1 $\approx$ 2{,}400, pc2 $\approx$ 1{,}150, pc3 $\approx$ 2{,}300,
pc4 $\approx$ 940--1{,}300 depending on stage.

\subsection{Degenerate (NaN) towers — pc4}
\label{sec:nan}

Two pc4 towers (cell IDs 26947081 and 28792072) produced
\texttt{refined\_val\_mae = NaN} in both Stage~3 and Stage~4. Diagnosis:
both have very small validation sets (23 and 9 rows respectively); at the
Nelder-Mead-refined position, Sionna found zero valid propagation paths
to \emph{every} row in that tower's validation subset (all predictions
hit the $-100$\,dBm sentinel floor), so even the code's flat-signal
fallback (mean of training RSRP) divides by an empty masked array and
returns NaN. This reproduced identically under scattering, confirming it
is a geometry/orientation degeneracy specific to these two
small-validation-set towers, not a scattering artifact. Both towers are
excluded from the Stage~3/4 aggregates above (45 candidate towers $\to$
43 valid). This is consistent with the paper's existing observation
(\S IV.E) that refinement most often misbehaves for towers with the
fewest training/validation rows.

\section{Gate 1: Full Statistics}
\label{sec:gate1full}

Supports: \S IV.A (Gate 1: Dead-Zone Preservation) in main.tex.

\subsection{Gap-event counts (gap-level, not row-level)}
\begin{center}
\begin{tabular}{lccc}
\toprule
Period & TypeA & TypeB & Ambiguous \\
\midrule
Train (days 1--2) & 4 & 57 & 13 \\
Val (day 3)        & 1 & 15 & 4 \\
\bottomrule
\end{tabular}
\end{center}
Training TypeA gaps split evenly by operator (2 from pc2/Op.\,2, 2 from
pc4/Op.\,1). The dominant TypeB cluster (56 of 57 training TypeB gaps)
belongs to Operator~2, centroid $(52.514212^\circ\text{N}, 13.348557^\circ\text{E})$,
stable to within 15\,m across all three collection days.

\subsection{Threshold sensitivity, full sweep}
\begin{center}
\begin{tabular}{lccc}
\toprule
$d_\text{recur}$ & Matched (of 15 val TypeB) & $C$ & Note \\
\midrule
50\,m  & 11/15 & 0.73 & strict; margin gaps excluded \\
100\,m (used) & 12/15 & 0.80 & midpoint, matches 1\,Hz GPS uncertainty \\
200\,m & 15/15 & 0.87\textsuperscript{a} & permissive \\
\bottomrule
\end{tabular}
\end{center}
\raggedright\footnotesize\textsuperscript{a}Reported as 0.87 in main text;
some sweep runs during development returned slightly different
denominators depending on whether ambiguous-adjacent gaps were
re-classified at wider radii — 0.87 is the value used in all published
figures.

\subsection{Multi-KPI corroboration, full statistics}
Held-out day 3, Operator~2, comparing rows within 150\,m of the TypeB
centroid (``near,'' 819 rows total / 723--745 with valid readings
depending on the KPI column) against the remaining route (``background,''
22{,}217--22{,}245 rows).

\begin{table}[h]
\centering
\begin{tabular}{lcccc}
\toprule
KPI & Near mean & Background mean & Welch $t$ & $p$ \\
\midrule
RSRP (dBm)     & $-115.43$ & $-89.97$ & $-64.22$  & $1.7\times10^{-316}$ \\
RSRQ (dB)      & $-15.99$  & $-10.48$ & $-32.87$  & $6.2\times10^{-147}$ \\
SNR$_1$ (dB)   & $-2.73$   & $10.60$  & $-42.25$  & $7.4\times10^{-202}$ \\
SNR$_2$ (dB)   & $-3.30$   & $9.40$   & $-53.58$  & $2.1\times10^{-269}$ \\
Avg.\ MCS index & $5.88$   & $15.08$  & $-60.18$  & $3.7\times10^{-315}$ \\
Datarate (bps) & $1.37\times10^6$ & $1.97\times10^7$ & $-134.30$ & $\approx 0$ \\
\bottomrule
\end{tabular}
\end{table}

\section{Twin-Gate Convergence: Pipeline History}
\label{sec:convergence_history}

Supports: \S IV.B (Twin-Gate Convergence) in main.tex. Reported here as a
robustness/consistency check: three successive, independently-computed
pipeline versions all land in a narrow band.

\begin{table}[h]
\centering
\small
\begin{tabular}{lcccc}
\toprule
Pipeline version & Predicted drop & Measured drop & \% recovered & Basis \\
\midrule
Early stratified-split & 27.2\,dB & 30.3\,dB & 90.0\% & pre-fix, both Op.\,2 devices \\
Intermediate (pc3-only) & 27.5\,dB & 31.5\,dB & 87.4\% & pc3 only, temporal, Stage 4 \\
Final (pc2+pc3) & 26.3\,dB & 30.4\,dB & 86.5\% & complete Op.\,2, temporal, Stage 4 \\
\bottomrule
\end{tabular}
\end{table}

Final version detail: near-zone mean predicted $-115.57$\,dBm, background
mean predicted $-89.29$\,dBm; near-zone mean measured $-120.02$\,dBm,
background mean measured $-89.65$\,dBm. 364 near-zone rows, 14{,}807
background rows (15{,}171 total, pc2+pc3 combined). Welch $t=-56.63$,
$p=6.2\times10^{-193}$.

\section{M2 Recompute: Provenance and Correction History}
\label{sec:m2history}

Supports: \S IV.C (M2 dead-zone detection) in main.tex.

\textbf{What happened.} The M2 figure originally reported (F1$=0.63$,
precision $0.52$, recall $0.82$, threshold $-116$\,dBm) was computed by
\texttt{M2\_serving\_cell.py} from \texttt{per\_tower\_val\_rows\_op1/op2.csv}
— files dated prior to this project's temporal-split correction, using
\emph{raw WCL positions with no Nelder-Mead refinement and no
scattering}. This was discovered while investigating why M2 seemed
disconnected from the corrected Gate 2 pipeline.

\textbf{Correction.} M2 was recomputed from the current Stage~4
(refined position + $S=0.4$ scattering) per-row predictions:

\begin{table}[h]
\centering
\begin{tabular}{lccccc}
\toprule
Basis & Threshold & Precision & Recall & F1 & TypeB rows \\
\midrule
Stale (WCL-only, pre-fix)      & $-116$\,dBm & 0.52 & 0.82 & 0.63 & 386/42{,}830 \\
pc3-only (Stage 4)             & $-113$\,dBm & 0.75 & 0.89 & 0.82 & 364/13{,}127 \\
pc2+pc3 combined (Stage 4, final) & $-113$\,dBm & 0.67 & 0.89 & 0.77 & 364/15{,}171 \\
\bottomrule
\end{tabular}
\end{table}

Note the pc3-only estimate (F1$=0.82$) is \emph{higher} than the combined
estimate (F1$=0.77$): pc2 contributes 2{,}044 additional background rows
but zero additional TypeB-zone rows on this held-out day (pc2's smaller,
differently-timed validation set did not pass through the dead zone),
so adding it can only add potential false positives, not true positives.
The combined figure is reported as the paper's headline M2 result because
it is the complete, honest measurement, not because it is the most
favorable one.

\section{Explored, Not Included: LOS/NLOS Classification and Multipath Richness (M4)}
\label{sec:m4}

\textbf{Status: real result, held out of main.tex by author decision, not
by any negative finding.} This section documents a fourth candidate
readiness axis — structural/mechanistic validation, distinct from M1
(accuracy) and M2 (localization) — computed via a new script (A17) that
extracts per-path richness metrics from Sionna's existing solver output
(delay, angle of arrival, interaction type) at zero additional RT cost,
reusing each device's already-refined Stage~4 geometry.

\subsection{Method}
Per validation row: power-weighted RMS delay spread ($\tau_\text{rms}$),
power-weighted circular RMS azimuth spread, valid path count, and an
LOS/NLOS flag (LOS if any valid path has an all-\texttt{NONE} interaction
sequence — i.e.\ a direct, unobstructed line — per Sionna's
\texttt{InteractionType} enum: \texttt{NONE}=0, \texttt{SPECULAR}=1,
\texttt{DIFFUSE}=2, \texttt{REFRACTION}=4, \texttt{DIFFRACTION}=8).

\subsection{M4a: LOS/NLOS accuracy}
\begin{table}[h]
\centering
\small
\begin{tabular}{lccccc}
\toprule
Device & LOS mean (dBm) & LOS std & NLOS mean (dBm) & NLOS std & Welch $t$ ($p$) \\
\midrule
pc1 & $-84.89$ & 10.83 & $-93.65$ & 9.07 & $25.6$ ($6\times10^{-109}$) \\
pc3 & $-90.12$ & 13.19 & $-81.93$ & 10.47 & $-24.2$ ($4\times10^{-109}$) \\
pc2 & $-96.47$ & 11.49 & $-95.53$ & 7.40 & $-1.4$ (n.s., $p=0.16$) \\
\bottomrule
\end{tabular}\\[2mm]
\raggedright\footnotesize
pc1: 81 towers, 14{,}172 rows. pc3: 63 towers, 13{,}127 rows. pc2: 36 towers, 2{,}044 rows.
\end{table}

\textbf{This is not a clean, consistent result across devices.} pc1 shows
LOS $>$ NLOS in measured RSRP, exactly as propagation physics predicts.
pc3 shows the \emph{opposite} sign, and pc2 shows no significant
difference at all. This inconsistency — not a single clean number — is
the actual reason this metric is not in main.tex: a metric that flips
sign between two devices on the same operator's network cannot be
presented as a validated readiness criterion without a lot more
investigation into why (candidate explanations, not yet tested:
Operator~2's more complex, NLOS-dominated urban geometry near the
confirmed dead zone; a sector-antenna-orientation interaction with which
paths get classified as ``direct''). \emph{Literature note}: VaN3Twin's
real-time follow-up (Pegurri et al., arXiv 2601.16559) does validate a
C-V2X DT using exactly this kind of LOS/NLOS transition-classification
metric, so this is a recognized axis in peer work, not an invented one —
worth revisiting with a cleaner methodology rather than abandoning.

\subsection{M4b: Multipath richness vs.\ M1 residual}
\begin{table}[h]
\centering
\begin{tabular}{lcccc}
\toprule
Device & corr($\tau_\text{rms}$, residual) & $p$ & corr($n_\text{valid}$, residual) & $p$ \\
\midrule
pc1 & $-0.032$ & $1.5\times10^{-4}$ & $-0.065$ & $6.3\times10^{-15}$ \\
pc3 & $-0.009$ & $0.308$ (n.s.)     & $-0.103$ & $4.6\times10^{-32}$ \\
pc2 & $-0.144$ & $5.9\times10^{-11}$ & $-0.143$ & $7.8\times10^{-11}$ \\
\bottomrule
\end{tabular}
\end{table}
All correlations are negligible in magnitude ($|r| < 0.15$) despite
significance at large $N$. \textbf{Null result}: multipath richness does
not meaningfully explain where M1 residual error concentrates. This
suggests the residual is dominated by something else (OLS calibration
slack, GPS positional noise, or building-material approximation error),
which is itself a useful negative finding pointing future work away from
richness-based residual explanations.

\section{Explored, Not Included: Handover / Cell-Boundary Prediction}
\label{sec:handover}

\textbf{Status: genuine negative result, well-diagnosed, not a bug.}
Motivated by the question ``does Gate 2's geometry correctly predict
\emph{where} a moving vehicle should hand over between cells,'' tested
two independent ways.

\subsection{Feasibility groundwork}
Gate~1's own TypeA (session-boundary) events are far too rare to serve as
ground truth: only 5 total across the entire 3-day dataset (4 train, 1
val) — the TypeA definition conflates handover with RLF-reestablishment
and logger restarts, and requires a measurement \emph{gap}, which most
real handovers do not cause. The better ground-truth signal is raw
serving-cell (\texttt{PCell\_Cell\_Identity}) transitions, filtered for
persistence to exclude cell-flapping noise:

\begin{center}
\begin{tabular}{lcccc}
\toprule
Persistence filter & pc1 val transitions & pc2 & pc3 & pc4 \\
\midrule
$\geq$5\,s  & 391 & 48 & 408 & 52 \\
$\geq$10\,s & 336 & 41 & 313 & 43 \\
$\geq$20\,s & 258 & 31 & 228 & 31 \\
\bottomrule
\end{tabular}
\end{center}
At the 20\,s threshold, 96.1\% of pc1's held-out transitions have both
the source and destination tower already covered by the existing
Stage~4 refined-position set — no new RT position optimization needed.

\subsection{A18: Cross-cell direction test}
Method: for each covered transition, evaluate Sionna's raw
(uncalibrated) predicted power from both the source and destination
tower at the last pre-transition GPS point and the first post-transition
point, using each tower's already-refined Stage~4 position (no
recalibration — every transition here is intra-operator, so towers share
a comparable EIRP/hardware family, which was the original rationale for
using raw power). Metric: Handover Direction Accuracy (HDA) — fraction
of transitions where the predicted-stronger tower matches the real
serving cell at \emph{both} points.

\textbf{Result (pc1, 10\,s filter, 322 covered transitions, 293 after
removing 29 sentinel-contaminated (`no path found') events):}
\begin{itemize}
  \item HDA (strict): 10.6\% (31/293)
  \item Ranking-flip rate (softer — does the margin move in the right
        direction, regardless of absolute correctness at either point):
        62.1\% (182/293)
  \item \texttt{correct\_A} (source tower wins \emph{before} the real
        switch): 44.0\% — \emph{below} chance
  \item \texttt{correct\_B} (destination tower wins \emph{after} the real
        switch): 65.2\% — above chance
\end{itemize}

\textbf{Diagnosis.} The asymmetry between \texttt{correct\_A} (below
chance) and \texttt{correct\_B} (above chance) is not random noise — it
points to a systematic effect. Each tower's position/height/azimuth was
refined by Nelder-Mead to minimize error over \emph{that tower's own
bulk of training rows}, which are concentrated in the interior of its
serving area. A cell-boundary point is, by construction, the sparsest,
least-represented part of each tower's own training distribution — the
edge where that tower is weakest and about to hand off. The
per-tower-independent calibration that gives a solid overall M1 (3.02\,dB
average) is systematically weakest exactly where handover-boundary
accuracy would need to be strong. This is a real limitation of the
current calibration architecture (independent per-tower fits), not a
bug in this experiment or evidence against the rest of the pipeline.

\subsection{Independent corroboration: pre-existing proxy\_b}
A simpler, independently-implemented handover check already existed in
\texttt{M2\_serving\_cell.py} (``Proxy B''), using \emph{calibrated}
(not raw) power at a single transition point, on the older stratified-split,
unrefined-position pipeline: \textbf{351 handover events, 156 correct,
44.4\% accuracy} — also below chance. Two independent implementations,
different calibration approaches, different pipeline versions, and
different data splits reach the same qualitative conclusion: cross-cell
handover-boundary prediction does not work well under
per-tower-independent Gate~2 calibration. This convergence of two
separate negative results is itself fairly strong evidence the limitation
is real and structural, not an implementation artifact — and settles the
question for the current architecture rather than inviting further quick
attempts at a fix.

\subsection{What would actually be needed}
A joint, boundary-aware calibration (e.g., fitting neighboring towers
together rather than independently, or explicitly up-weighting training
rows near candidate boundaries) — a genuinely new methodology, not a
parameter tweak, and a reasonable candidate for a journal extension
rather than this conference submission.

\section{Literature Scan Summary}
\label{sec:litscan}

A dedicated literature scan (\texttt{DT\_READINESS\_LITERATURE\_SCAN.md},
in the project repository) surveyed published Digital Twin
implementations for C-V2X / vehicular cellular coverage to identify
which experiments and thresholds TWINGATE should consider replicating.
Key findings:

\begin{itemize}
  \item \textbf{DT-CoVeSS} (Twardokus \& Rahbari, published IEEE INFOCOM
        Workshops/DTwin 2025) is a Sionna RT DT built specifically for
        \emph{PC5 sidelink} — confirms sidelink C-V2X DT validation is a
        distinct research track with its own data requirements,
        independently supporting TWINGATE's own scope decision to stay
        Uu-mode only given the Berlin V2X dataset has no sidelink logs.
  \item \textbf{Lucas-Estañ et al.} (published IEEE \emph{Network}, vol.
        37, no.\ 4, 2023) grounds V2X service-continuity requirements in
        exact 3GPP TS~22.186 numbers: Cooperative Lane Change requires
        90\% of packets under 25\,ms (low automation) or 99.99\% under
        10\,ms (high automation). This is the strongest available
        motivation for the handover/cell-boundary direction
        (\S\ref{sec:handover}), even though that direction did not pan
        out under the current architecture.
  \item \textbf{Pegurri et al.} (VaN3Twin real-time follow-up, arXiv
        2601.16559) validates LOS/NLOS transition-classification
        accuracy in a live 60\,GHz C-V2X testbed — direct peer precedent
        for the M4 direction (\S\ref{sec:m4}).
  \item No paper found translates a 3GPP V2X reliability/latency target
        into an RT-usable RSRP/SINR coverage threshold; doing so would
        require a link-level BLER-vs-SINR simulation step (a different
        Sionna module — PHY/LDPC, not RT), which is a materially larger
        undertaking flagged as a longer-horizon journal-extension item.
\end{itemize}

\section{Full Limitations and Future Work Detail}
\label{sec:limitations_full}

Supports: \S V (Discussion, Limitations subsection) in main.tex, which
compresses these to one-line bullets for space. Full reasoning below.

\subsection{Limitations, expanded}

\textbf{Single-dataset scope.} All results are from one urban route in
Berlin (17.2\,km loop, two LTE operators, June 2021). The pipeline
design is dataset-agnostic by construction, but we make no empirical
claim about performance on other cities, network generations (5G NR),
or high-mobility scenarios (highway V2X at $>$100\,km/h). The TypeB
coverage hole identified by Gate~1 is specific to the Vodafone network
in this deployment; other datasets may exhibit different gap
distributions or none at all.

\textbf{LOS/NLOS path decomposition and dynamic scatterers.} Sionna~RT's
path solver natively classifies every propagation path by interaction
type (line-of-sight, specular reflection, diffraction), so
LOS/NLOS-conditioned RSRP statistics are directly extractable from the
existing pipeline without architectural changes (see \S\ref{sec:m4} for
what was found when this was actually tried). Moving vehicles in a
dense platoon act as both secondary blockers and reflectors; because the
same GPS segments are traversed by multiple devices across multiple
days, instantaneous blockage events are partially de-correlated in the
per-segment RSRP distributions used for Gate~2 calibration. A tractable
extension is to add each companion vehicle as a dynamic mesh object at
its GPS-synchronized timestamp, converting the batched PathSolver call
to per-row ray tracing; the synchronised multi-device tracks in this
dataset make this feasible, at increased per-tower compute cost (see
\S\ref{sec:handover} for the concrete feasibility numbers worked out
for a related dynamic-scene idea this session).

\textbf{Wideband metrics and Doppler.} At 1.8\,GHz LTE and a vehicle
speed of 100\,km/h, the maximum Doppler shift is approximately 167\,Hz,
giving a channel coherence time of roughly 3\,ms. The 3GPP TS~36.133 L3
measurement filter averages RSRP over a 200\,ms window, spanning on the
order of 60 coherence times: fast Doppler fading therefore averages out
in the RSRP metric validated against, and explicit Doppler modeling is
not required for coverage-level DT readiness assessment. Wideband
metrics (delay spread, channel Doppler spectrum, block-error rate) do
require dynamic scatterer modeling; MobileInsight does not log I/Q
samples, placing these metrics outside the scope of the present
framework.

\textbf{Search radius.} The Gate~2 optimizer uses a $\pm$500\,m
horizontal search radius centered on the WCL crowd-sourced
initialization. Towers serving peripheral route segments, where
drive-test GPS coverage is sparse and WCL positions are correspondingly
less reliable, may fall outside this radius and receive sub-optimal
localization. Widening the radius increases the risk of convergence to
distant spurious positions; a two-stage coarse-to-fine strategy would
address this trade-off.

\textbf{Computational cost.} The Nelder-Mead optimization loop in
Gate~2 requires 3--6 minutes per tower on a single CPU core, totaling
several hours for the full multi-hundred-tower deployment. This is
acceptable for offline DT construction but unsuitable for real-time
update loops or rapid re-deployment scenarios. Because each tower's
optimization is fully independent, wall time scales linearly with
available cores; learned surrogate models for the Sionna forward pass
offer a complementary path to per-tower cost reduction.

\subsection{Future work, expanded}

\textbf{Dynamic scene extension.} The most direct path to capturing
vehicular blockage and reflection is to overlay a stochastic
mobile-obstacle model on the static Sionna~RT scene. Each vehicle in
the platoon can be modeled as a rectangular lossy-dielectric slab whose
position along the route is drawn from a Poisson point process
calibrated to observed traffic density; blockage probability as a
function of inter-vehicle separation follows from Poisson thinning and
can be computed analytically, adding a per-link stochastic attenuation
term without modifying the ray-tracing engine. (Note: this session also
scoped and ruled out a related but distinct idea -- real companion-vehicle
mesh objects driven by actual GPS tracks rather than a statistical
model -- see \S\ref{sec:handover}'s feasibility groundwork; the
statistical overlay proposed here avoids that approach's per-row
scene-reload cost.)

\textbf{M3: a channel-dynamics readiness metric.} A third readiness
metric based on Doppler spread would shift validation from signal-level
to channel-level fidelity. Doppler spread can in principle be estimated
from per-TTI RSRQ fluctuations already logged by MobileInsight, making
M3 accessible without requiring I/Q samples. (Note: this specific idea
was investigated this session and found not currently viable with this
dataset -- the cleaned CSV is uniformly 1\,Hz, three orders of magnitude
too coarse to resolve millisecond-scale coherence times. A genuine M3
would need either raw sub-second MobileInsight logs, if recoverable, or
a different dataset.)

\textbf{Online calibration.} The 3--6\,min/tower offline cost of Gate~2
can be amortized into a continuous online loop using recursive
least-squares updates on the per-tower OLS offsets: as vehicles traverse
the route, incoming RSRP measurements refine each tower's calibration
incrementally, converging to a solution equivalent to full
re-optimization within a few vehicle passes and eliminating the need
for periodic batch re-runs.

\section{Miscellaneous Fixes Made During This Project}
\label{sec:fixes}

For completeness/transparency, corrections made to the pipeline and
paper during development, beyond the M2/convergence recompute already
covered in \S\ref{sec:m2history}:

\begin{itemize}
  \item \textbf{Dataset year}: a transcription error in early drafts
        referred to the Berlin V2X collection dates as ``June 2022'';
        corrected to the true ``22--24 June 2021'' throughout.
  \item \textbf{Gate 1 completeness $C$}: an early draft reported
        $C=0.70$ using an incorrect denominator (all validation gaps,
        including TypeA/Ambiguous); corrected to $C=0.80$ using the
        TypeB-only denominator specified in main.tex's completeness
        equation (Section~III-B.3).
  \item \textbf{Cosmetic logging bugs}: \texttt{A15\_tx\_position\_refinement.py}
        had two hardcoded string labels (``Op1'' in the scene-load
        message, ``pc1'' in the final-results banner) that did not
        reflect the actual device/operator being processed — purely
        cosmetic (the underlying \texttt{SCENE\_XML} path and output
        filenames were always correct), fixed to use the actual
        \texttt{DEVICE}/\texttt{OPERATOR} variables.
\end{itemize}

\end{document}
